% Section 7: Conclusion (~600 words)

\subsection{Summary}

This paper synthesized five empirical validation studies establishing
the robustness of METIS recursive bisection redistricting across four
dimensions. The key quantitative findings are:

\begin{itemize}
  \item \textbf{Spatial robustness (C.1)}: Mean $\Delta PP < 0.005$
    across a 130$\times$ unit-count range. Geography drives outcomes,
    not resolution.

  \item \textbf{Temporal stability (C.2)}: Geographic variance is
    $3.2\times$ larger than temporal variance. Within-slice cross-census
    correlation is $r = 0.73$. National PP varies by $\approx 10\%$
    over 20 years.

  \item \textbf{Boundary stability (C.3)}: Recursive bisection provides
    $0.8$ pp better decadal stability than n-way partitioning. IoU
    across census cycles is $71\%$ within geographic slices.

  \item \textbf{Longitudinal convergence (C.4)}: The algorithmic-enacted
    gap narrowed from $0.114$ (2000) to $0.110$ (2020) as reform
    spreads. Geographic clustering of algorithm performance is stable:
    47/50 states remain in the same cluster across 20 years.

  \item \textbf{Partisan fairness (C.5)}: Mean $|EG| = 0.04$ for
    algorithmic plans vs.\ $0.08$ for enacted. No algorithmic plan
    exceeds the $7\%$ court threshold; 8 of 15 enacted plans do.
    Five independent metrics converge on an $\approx 8$ pp partisan
    gap.
\end{itemize}

\subsection{The Validation Portfolio as Legal Evidence}

The C-track is designed to be cited as a unit in legal proceedings
where algorithmic redistricting is challenged or defended. A challenge
that the algorithm produces different results at different resolutions
is answered by C.1. A challenge that it fails to generalize across
census years is answered by C.2 and C.3. A challenge that it is
implicitly partisan is answered by C.5. A challenge that its
performance clusters by geography in ways that disadvantage particular
states is answered by C.4.

No single paper is sufficient; the combination provides a
multi-dimensional defense that matches the multi-dimensional nature
of redistricting litigation. Courts examining algorithmic redistricting
claims under state constitutions require evidence on compactness,
population balance, VRA compliance, and partisan fairness. The C-track
provides quantitative baselines for all four.

\subsection{The Broader Point: Validation Is Not Optional}

The DIA's certification requirements reflect a deeper truth: without
systematic validation, algorithmic redistricting is an assertion, not
a proof. An algorithm that has been run once on one census year and
found to produce compact districts is a promising tool; an algorithm
that has been systematically tested across three spatial resolutions,
three census years, five fairness metrics, and found robust on all of
them is a validated tool appropriate for use in consequential legal
proceedings.

The C-track constitutes, to our knowledge, the first systematic
multi-dimensional validation program for an algorithmic redistricting
method at national scale. We hope it serves as a template for other
researchers developing redistricting algorithms: robustness claims
require robustness evidence, and robustness evidence requires the kind
of systematic, multi-dimensional empirical testing this track exemplifies.

\subsection{Future Work}

Three extensions to the C-track are planned.

First, completing block-level analysis (C.1) for all 50 states will
close the 10-state gap in spatial robustness validation. This requires
approximately 50 additional compute-hours using the Rust \texttt{redist}
binary.

Second, incorporating 2030 Census data into C.3 and C.4 will provide
a third temporal data point, enabling trend analysis rather than
level-shift characterization.

Third, extending C.5's partisan fairness analysis to all 50 states
(including small states with few districts) using multi-level modeling
will complete the national partisan fairness picture. Preliminary
analysis suggests small states' single-district or two-district plans
exhibit near-zero EG by construction; the key variation is in
medium-sized states (5--15 districts) where algorithmic and enacted
differences are largest.
